\subsection{\texorpdfstring{Hilbert 空间 \(L^{2}(G)\) 的情形}{Hilbert 空间 L\^{}\{2\}(G) 的情形}}

对 \(\widehat{\mathbf{G}}\) 的每个可测子集 M，用 \(\mathrm{E_M}\) 表示 \(f \in \mathrm{L}^2(\mathbf{G})\)（其 Fourier 变换 \(\mathcal{F}_{\mathrm{G}}(f)\) 在 \(\widehat{\mathbf{G}}\) 上几乎处处为零）所成的集合。设 \(\varphi_{\mathrm{M}}\) 是 M 的特征函数。空间 \(\mathrm{E_M}\) 是连续线性映射 \(f \mapsto \varphi_{\mathrm{M}}\mathcal{F}_{\mathrm{G}}(f)\)（由 \(\mathrm{L}^2(\mathbf{G})\) 到 \(\mathrm{L}^2(\widehat{\mathbf{G}})\) 中）的核，因而它是 \(\mathrm{L}^2(\mathbf{G})\) 的一个闭子空间。

命题 1.　a) 设 M 是 \(\widehat{\mathbf{G}}\) 的一个可测子集。对每个 \(x \in \mathbf{G}\)，空间 \(\mathrm{E}_{\mathrm{M}}\) 在映射 \(f \mapsto \varepsilon_x * f\) 之下稳定；

\begin{enumerate}
\def\labelenumi{\alph{enumi})}
\setcounter{enumi}{1}
\item
  设 M 与 N 是 G 的可测子集。有 \(E_{M} = E_{N}\)，当且仅当 M 与 N 在相差一个局部可忽略集的意义下相等；
\item
  \(L^{2}(G)\) 的每个在诸映射 \(f \mapsto \varepsilon_{x} * f\)（对所有 \(x \in G\)）之下稳定的子空间都具有 \(E_{M}\) 的形式，其中 M 是 \(\widehat{G}\) 的某个可测子集。
\end{enumerate}

这一结果将在以后证明（参看 V，待出版）。
% semantic-fragments: legacy-input-seam
\relax{}
